\documentclass[10pt]{article}
\usepackage[T1]{fontenc}
\usepackage[utf8]{inputenc}
\usepackage[french]{babel}
\usepackage{amsmath,amssymb,mathrsfs,mathpazo}
\usepackage{geometry}
\geometry{a4paper,margin=2cm}
\usepackage{graphicx,array,tabularx,tabularray,colortbl,multirow,multicol,booktabs,enumitem,pifont,ulem,fancybox,diagbox,cancel}
\usepackage{pstricks,pstricks-add,pst-plot,pst-tree,pst-node,pst-text,pst-func,pst-math,pst-bspline,pst-3dplot,pst-eucl,pst-solides3d}
\let\pscalculate\undefined
\usepackage{pst-bezier}
\usepackage{tikz,pgfplots,tkz-tab}
\usetikzlibrary{arrows,arrows.meta,patterns,calc,positioning,trees,decorations.markings,shapes.geometric,shapes.misc,shapes.arrows,intersections,decorations.pathmorphing,decorations.pathreplacing,backgrounds,shadows,babel}
\pgfplotsset{compat=1.18}
\usepgfplotslibrary{fillbetween}
\usepackage[np]{numprint}
\usepackage{listings,makecell,pas-tableur,twoopt}
\let\barre\undefined
\usepackage[pstricks]{tablvar}
\definecolor{darkgreen}{RGB}{0,100,0}
\providecommand{\red}{\color{red}}
\providecommand{\blue}{\color{blue}}
\providecommand{\green}{\color{green}}
\providecommand{\black}{\color{black}}
\providecommand{\white}{\color{white}}
\providecommand{\gray}{\color{gray}}
\providecommand{\cyan}{\color{cyan}}
\providecommand{\magenta}{\color{magenta}}
\providecommand{\yellow}{\color{yellow}}
\providecommand{\violet}{\color{violet}}
\providecommand{\orange}{\color{orange}}
\providecommand{\euro}{\texteuro}
\providecommand{\EUR}{\texteuro}
\providecommand{\ieme}{\textsuperscript{e}}
\providecommand{\degres}{\ensuremath{^\circ}}
\providecommand{\degree}{\ensuremath{^\circ}}
\providecommand{\vv}[1]{\overrightarrow{#1}}
\providecommand{\Coord}[1]{#1}
\providecommand{\decosix}{}
\providecommand{\starredbullet}{\ensuremath{\bullet}}
\pagestyle{empty}
\setlength{\parindent}{0pt}
\usepackage{adjustbox}\DeclareUnicodeCharacter{2236}{:}\DeclareUnicodeCharacter{2190}{\ensuremath{\leftarrow}}
\providecommand{\R}{\mathbb{R}}
\providecommand{\N}{\mathbb{N}}
\providecommand{\Z}{\mathbb{Z}}
\providecommand{\Q}{\mathbb{Q}}
\providecommand{\D}{\mathbb{D}}
\providecommand{\C}{\mathbb{C}}
\providecommand{\e}{{\mathrm{e}}}
\providecommand{\vect}[1]{\overrightarrow{#1}}
\providecommand{\vectt}[1]{\overrightarrow{\mathrm{#1}}}
\providecommand{\vvt}[1]{\overrightarrow{\mathrm{#1}}}
\providecommand{\np}[2][]{{#2~#1}}
\providecommand{\vv}[1]{{\overrightarrow{#1}}}
\renewcommand{\d}{{\mathrm{d}}}
\providecommand{\degres}{{^\circ}}
\providecommand{\degre}{{^\circ}}
\providecommand{\blue}{\color{blue}}
\providecommand{\red}{\color{red}}
\providecommand{\enskip}{\quad}
\providecommand{\llbracket}{\lbrack\!\lbrack}
\providecommand{\rrbracket}{\rbrack\!\rbrack}
\providecommand{\euro}{{\text{€}}}
\providecommand{\no}{{\text{n°}}}
\providecommand{\interval}[2]{{[#1;#2]}}
\providecommand{\intervFO}[2]{{[#1;#2[}}
\providecommand{\dotfill}{\ldots}

\providecommand{\ds}{\displaystyle}
\providecommand{\textding}[1]{#1}
\providecommand{\Oij}{\(\mathrm{O};\vec{\imath},\vec{\jmath}\)}
\providecommand{\Oijk}{\(\mathrm{O};\vec{\imath},\vec{\jmath},\vec{k}\)}
\providecommand{\Ouv}{\(\mathrm{O};\vec{u},\vec{v}\)}
\providecommand{\intervFF}[2]{\ensuremath{[#1;#2]}}
\renewcommand{\euro}{\ifmmode\text{\texteuro}\else\texteuro\fi}
\DeclareUnicodeCharacter{00B0}{\ifmmode^{\circ}\else\textdegree\fi}
\DeclareUnicodeCharacter{20AC}{\ifmmode\text{\texteuro}\else\texteuro\fi}
% Énoncés : Éducation nationale. Transcription et corrigés : APMEP, Denis Vergès et contributeurs. Découpage et classement : CoopMaths, Sébastien Lozano, Éric Elter et contributeurs. Adaptation de la présentation : Maths Explorer.
% CC BY-SA : https://coopmaths.fr/www/about/
% https://coopmaths.fr/alea/static/bac/2025/tex/bac_2025_09_sujet1_asie_4.tex
\begin{document}
\noindent\textit{APMEP / CoopMaths. CC BY-SA. Présentation adaptée par Maths Explorer.}\par\medskip
Soit $n$ un entier naturel non nul.

Dans le cadre d'une expérience aléatoire, on considère une suite d'évènements $A_{n}$ et on note $p_{n}$ la probabilité de l'évènement $A_{n}$.

Pour les parties \textbf{A} et \textbf{B} de l'exercice, on considère que :

\begin{itemize}
\item Si l'évènement $A_{n}$ est réalisé alors l'évènement $A_{n+1}$ est réalisé avec une probabilité 0,3.
\item Si l'évènement $A_{n}$ n'est pas réalisé alors l'évènement $A_{n+1}$ est réalisé avec une probabilité 0,7.
\end{itemize}

On suppose que $p_{1}=1$.

\subsection*{Partie A :}

\begin{enumerate}
\item Recopier et compléter les probabilités sur les branches de l'arbre des probabilités ci-dessous:

\begin{center}
	\begin{tikzpicture}[xscale=1,yscale=1,baseline={(R.base)}]
	
	\tikzstyle{fleche}=[->,>=latex,thick]
	\tikzstyle{noeud}=[fill=white,circle,inner sep=2pt]
	\tikzstyle{feuille}=[fill=white,circle,inner sep=2pt]
	\tikzstyle{etiquette}=[pos=0.6,fill=white, inner xsep=3pt, inner ysep=1.5pt]
	
	\def\DistanceInterNiveaux{3}
	\def\DistanceInterFeuilles{0.8}
	
	\def\NiveauA{(0)*\DistanceInterNiveaux}
	\def\NiveauB{(1)*\DistanceInterNiveaux}
	\def\NiveauC{(2)*\DistanceInterNiveaux}
	\def\InterFeuilles{(-1)*\DistanceInterFeuilles}
	
	\node[noeud] (R) at ({\NiveauA},{(1.5)*\InterFeuilles}) {$A_1$};
	\node[noeud] (Ra) at ({\NiveauB},{(0.5)*\InterFeuilles}) {$A_2$};
	\node[feuille] (Raa) at ({\NiveauC},{(0)*\InterFeuilles}) {$A_3$};
	\node[feuille] (Rab) at ({\NiveauC},{(1)*\InterFeuilles}) {$\overline{A_3}$};
	\node[noeud] (Rb) at ({\NiveauB},{(2.5)*\InterFeuilles}) {$\overline{A_2}$};
	\node[feuille] (Rba) at ({\NiveauC},{(2)*\InterFeuilles}) {$A_3$};
	\node[feuille] (Rbb) at ({\NiveauC},{(3)*\InterFeuilles}) {$\overline{A_3}$};
	
	\draw[fleche] (R.east)--(Ra.west);
	\draw[fleche] (Ra.east)--(Raa.west);
	\draw[fleche] (Ra.east)--(Rab.west);
	\draw[fleche] (R.east)--(Rb.west) node[etiquette] {$0,7$};
	\draw[fleche] (Rb.east)--(Rba.west);
	\draw[fleche] (Rb.east)--(Rbb.west);
	\end{tikzpicture}
\end{center}

\item Montrer que $p_{3}=0,58$.

\item Calculer la probabilité conditionnelle $P_{A_{3}}\left(A_{2}\right)$, arrondir le résultat à $10^{-2}$ près.
\end{enumerate}

\subsection*{Partie B :}

	Dans cette partie, on étudie la suite $\big(p_{n}\big)$ avec $n \geqslant 1$.

\begin{enumerate}
\item Recopier et compléter les probabilités sur les branches de l'arbre des probabilités ci-dessous :

\begin{center}
		\begin{tikzpicture}[xscale=1,yscale=1,baseline={(R.base)}]
			
			\tikzstyle{fleche}=[->,>=latex,thick]
			\tikzstyle{noeud}=[fill=white,circle,inner sep=2pt]
			\tikzstyle{feuille}=[fill=white,circle,inner sep=2pt]
			\tikzstyle{etiquette}=[pos=0.6,fill=white, inner xsep=3pt, inner ysep=1.5pt]
			
			\def\DistanceInterNiveaux{3}
			\def\DistanceInterFeuilles{0.8}
			
			\def\NiveauA{(0)*\DistanceInterNiveaux}
			\def\NiveauB{(1)*\DistanceInterNiveaux}
			\def\NiveauC{(2)*\DistanceInterNiveaux}
			\def\InterFeuilles{(-1)*\DistanceInterFeuilles}
			
			\node[noeud] (R) at ({\NiveauA},{(1.5)*\InterFeuilles}) {};
			\node[noeud] (Ra) at ({\NiveauB},{(0.5)*\InterFeuilles}) {$A_n$};
			\node[feuille] (Raa) at ({\NiveauC},{(0)*\InterFeuilles}) {$A_{n+1}$};
			\node[feuille] (Rab) at ({\NiveauC},{(1)*\InterFeuilles}) {$\overline{A_{n+1}}$};
			\node[noeud] (Rb) at ({\NiveauB},{(2.5)*\InterFeuilles}) {$\overline{A_n}$};
			\node[feuille] (Rba) at ({\NiveauC},{(2)*\InterFeuilles}) {$A_{n+1}$};
			\node[feuille] (Rbb) at ({\NiveauC},{(3)*\InterFeuilles}) {$\overline{A_{n+1}}$};
			
			\draw[fleche] (R.east)--(Ra.west) node[etiquette] {$p_n$};
			\draw[fleche] (Ra.east)--(Raa.west);
			\draw[fleche] (Ra.east)--(Rab.west);
			\draw[fleche] (R.east)--(Rb.west) ;
			\draw[fleche] (Rb.east)--(Rba.west);
			\draw[fleche] (Rb.east)--(Rbb.west);
		\end{tikzpicture}
	\end{center}
	\item \begin{enumerate}
		\item Montrer que, pour tout entier naturel $n$ non nul : $p_{n+1}=-0,4 p_{n}+0,7$.
		\end{enumerate}
	On considère la suite $\big(u_{n}\big)$, définie pour tout entier naturel $n$ non nul par : 

$u_{n}=p_{n}-0,5$.

	\begin{enumerate}[resume]
	\item Montrer que $\big(u_{n}\big)$ est une suite géométrique dont on précisera la raison et le premier terme.

	\item En déduire l'expression de $u_{n}$, puis de $p_{n}$ en fonction de $n$.

	\item Déterminer la limite de la suite $\big(p_{n}\big)$.
	\end{enumerate}
\end{enumerate}

\subsection*{Partie C :}

Soit $x \in ]0 ~;~ 1[$, on suppose que $P_{\overline{A_{n}}}\left(A_{n+1}\right)=P_{A_{n}}\left(\overline{A_{n+1}}\right)=x$. On rappelle que $p_{1}=1$.

\begin{enumerate}
\item Montrer que pour tout entier naturel $n$ non nul : $p_{n+1}=(1 - 2x) p_{n}+x$.

\item Démontrer par récurrence sur $n$ que, pour tout entier naturel $n$ non nul :
		\[p_{n}=\dfrac{1}{2}(1-2 x)^{n-1}+\dfrac{1}{2}\]

\item Montrer que la suite $\big(p_{n}\big)$ est convergente et donner sa limite.
\end{enumerate}

\end{document}
